\section{Introduction}
\label{sec:introduction}

MQTT is the dominant messaging protocol for Internet of Things (IoT) deployments, widely adopted for device-to-cloud communication across industrial, smart city, and consumer IoT platforms~\cite{bansal2020iot-mqtt-survey}. MQTT's publish-subscribe model, lightweight framing, and multiple quality-of-service levels make it well-suited for constrained devices and unreliable networks~\cite{mqtt-v5-spec}.

However, MQTT operates over TCP, which introduces a fundamental limitation: \emph{head-of-line (HOL) blocking}~\cite{rfc9000, kumar2019mqtt-quic}. When an IoT device publishes to multiple topics over a single TCP connection---as is common for devices reporting temperature, humidity, GPS, and status simultaneously---a packet loss on any one topic's data blocks delivery of \emph{all} topics until retransmission completes. This creates artificial coupling between logically independent data flows, increasing tail latency and reducing the responsiveness of multi-topic workloads.

QUIC~\cite{rfc9000}, the transport protocol underlying HTTP/3, addresses HOL blocking through \emph{stream multiplexing}: multiple independent streams share a single connection, each with its own flow control and loss recovery. A lost packet on one stream does not block data on other streams. This property has been extensively studied for web traffic~\cite{langley2017quic, kakhki2017quic, marx2020hol, yu2021quic}, but its application to IoT messaging protocols remains largely unexplored.

Running MQTT over QUIC is not simply a matter of replacing the TCP socket. The question of \emph{how} to map MQTT's control and data packets onto QUIC streams admits multiple answers, each representing a different point in the isolation--overhead tradeoff space (Section~\ref{sec:architecture}).

Our contributions are:
\begin{enumerate}
\item Three stream mapping strategies (control-only, per-topic, per-publish) for MQTT over QUIC with a lightweight flow header protocol for stream demultiplexing.
\item The first empirical evaluation of MQTT stream mapping strategies under controlled network impairment, quantifying HOL blocking, throughput, and resource overhead.
\item An evaluation of QUIC datagram transport for MQTT as an alternative to stream-based delivery.
\item An open-source implementation with integrated benchmarking tools.
\end{enumerate}

The remainder of this paper is organized as follows. Section~\ref{sec:background} covers MQTT, QUIC, and related work. Section~\ref{sec:architecture} describes the stream mapping architecture. Section~\ref{sec:implementation} details the implementation. Section~\ref{sec:evaluation} presents experimental results. Section~\ref{sec:discussion} discusses strategy selection guidelines and limitations. Section~\ref{sec:conclusion} concludes.
